\documentclass[a4paper,10pt]{article}
\usepackage{lamia-pedagogy}

\begin{document}

% =========================================================================
% PAGE 1 : Fiche Pédagogique - Modèle Standard
% =========================================================================
\ficheHeader{1 Heure}{\makebox[2.5cm]{\dotfill}}{1\textsuperscript{ère} AM}{2026 / 2027}

\vspace{-3pt}
\competenceBanner{\faBookOpen\ Domaine : Grammaire / Activité de Langue}

\vspace{-2pt}
\modaliteSupportBar{Travail individuel / en binôme}{Manuel scolaire p. 10 / Cahier d'activités}

\vspace{1pt}

\begin{minipage}[t]{0.48\textwidth}
\begin{tcolorbox}[
  enhanced,
  colback=secondarylight,
  colframe=secondary,
  arc=1.5mm,leftrule=3pt,rightrule=0.5pt,toprule=0.5pt,bottomrule=0.5pt,
  top=3pt,bottom=3pt,left=6pt,right=6pt,
  title={\small\textbf{\faBullseye\ Objectifs d'apprentissage}},
  coltitle=secondary!90!black,attach boxed title to top left={yshift=-2mm,xshift=4mm},
  boxed title style={colback=white,boxrule=0.5pt,colframe=secondary,arc=1mm,top=1pt,bottom=1pt}
]
\vspace{1pt}
{\small
\begin{itemize}[leftmargin=*,noitemsep,topsep=1pt]
  \item[\faCheckCircle] Identifier la notion étudiée dans un texte court.
  \item[\faCheckCircle] Distinguer les formes et les règles d'accord.
  \item[\faCheckCircle] Réemployer correctement la notion dans des phrases.
\end{itemize}
}
\end{tcolorbox}
\end{minipage}%
\hfill
\begin{minipage}[t]{0.48\textwidth}
\begin{tcolorbox}[
  enhanced,
  colback=accentlight!35,
  colframe=accent,
  arc=1.5mm,leftrule=3pt,rightrule=0.5pt,toprule=0.5pt,bottomrule=0.5pt,
  top=3pt,bottom=3pt,left=6pt,right=6pt,
  title={\small\textbf{\faLightbulb\ Mémo / Règle d'or}},
  coltitle=accent!90!black,attach boxed title to top left={yshift=-2mm,xshift=4mm},
  boxed title style={colback=white,boxrule=0.5pt,colframe=accent,arc=1mm,top=1pt,bottom=1pt}
]
\vspace{1pt}
{\small
\textbf{Je retiens :}
\begin{itemize}[leftmargin=*,noitemsep,topsep=1pt]
  \item Règle essentielle claire et synthétique.
  \item Exemple clé : \textit{Le lion est un animal sauvage.}
  \item Attention aux pièges courants et confusions.
\end{itemize}
}
\end{tcolorbox}
\end{minipage}

\vspace{4pt}

% Section Activité 1
\begin{tcolorbox}[
  enhanced,
  colback=white,
  colframe=cardborder,
  arc=2mm,boxrule=0.8pt,
  top=4pt,bottom=4pt,left=7pt,right=7pt,
  title={\small\textbf{\textcolor{primary}{\faPen\ Activité 01 : Observation et repérage}} \hfill \badge[primary]{Compréhension}},
  coltitle=primary,colbacktitle=primarylight,
  titlerule=0.5pt,titlerule style=cardborder
]
{\small
\textbf{Consigne :} Lis les phrases suivantes puis réponds aux questions.
\begin{enumerate}[label=\textbf{\arabic*.},leftmargin=*,noitemsep,topsep=2pt]
  \item Exemple d'énoncé avec une question à choix multiples :
    \cardOptThree{Option A}{Option B}{Option C}
  \item Deuxième énoncé demandant un repérage précis :
    \cardOptTwo{Vrai}{Faux}
\end{enumerate}
}
\end{tcolorbox}

\vspace{3pt}

% Section Activité 2
\begin{tcolorbox}[
  enhanced,
  colback=white,
  colframe=cardborder,
  arc=2mm,boxrule=0.8pt,
  top=4pt,bottom=4pt,left=7pt,right=7pt,
  title={\small\textbf{\textcolor{primary}{\faEdit\ Activité 02 : Application dirigée}} \hfill \badge[emerald]{Entraînement}},
  coltitle=primary,colbacktitle=primarylight,
  titlerule=0.5pt,titlerule style=cardborder
]
{\small
\textbf{Consigne :} Complète les pointillés par le mot ou la forme qui convient.
\begin{enumerate}[label=\textbf{\alph*)},leftmargin=*,noitemsep,topsep=2pt]
  \item La phrase modèle à compléter : \drule{3.5cm} au bon endroit.
  \item Deuxième exemple pour l'élève : \drule{3.5cm} selon la règle.
  \item Troisième phrase avec réponse attendue : \drule{3.5cm}.
\end{enumerate}
}
\end{tcolorbox}

\vfill
\cutLineH

\newpage

% =========================================================================
% PAGE 2 : Fiche Pédagogique - Réinvestissement & Évaluation
% =========================================================================
\ficheEleveHeader

\vspace{-3pt}
\competenceBanner{\faPenFancy\ Réinvestissement et Production Écrite}

\vspace{2pt}

% Section Activité 3
\begin{tcolorbox}[
  enhanced,
  colback=white,
  colframe=cardborder,
  arc=2mm,boxrule=0.8pt,
  top=5pt,bottom=5pt,left=8pt,right=8pt,
  title={\small\textbf{\textcolor{primary}{\faLayerGroup\ Activité 03 : Consolidation}} \hfill \badge[secondary]{Approfondissement}},
  coltitle=primary,colbacktitle=primarylight,
  titlerule=0.5pt,titlerule style=cardborder
]
{\small
\textbf{Consigne :} Relie chaque élément de la colonne A à son correspondant dans la colonne B.

\vspace{3pt}
\begin{center}
\begin{tabularx}{0.92\linewidth}{|X|c|X|}
\hline
\rowcolor{primarylight}\textbf{Colonne A (Sujet)} & $\longleftrightarrow$ & \textbf{Colonne B (Prédicat / Verbe)} \\
\hline
1. Les élèves de la classe & $\bullet\quad\bullet$ & a. ont apporté leurs cahiers. \\
\hline
2. Le professeur de français & $\bullet\quad\bullet$ & b. explique la nouvelle leçon. \\
\hline
3. Nous & $\bullet\quad\bullet$ & c. sommes attentifs aux consignes. \\
\hline
\end{tabularx}
\end{center}
}
\end{tcolorbox}

\vspace{4pt}

% Section Activité 4
\begin{tcolorbox}[
  enhanced,
  colback=white,
  colframe=cardborder,
  arc=2mm,boxrule=0.8pt,
  top=5pt,bottom=5pt,left=8pt,right=8pt,
  title={\small\textbf{\textcolor{primary}{\faPenSquare\ Activité 04 : Situation d'intégration}} \hfill \badge[accent]{Production}},
  coltitle=primary,colbacktitle=primarylight,
  titlerule=0.5pt,titlerule style=cardborder
]
{\small
\textbf{Consigne :} Rédige 2 à 3 phrases pour présenter ton animal préféré en utilisant les notions vues aujourd'hui.

\vspace{6pt}
1. \dotfill\\[10pt]
2. \dotfill\\[10pt]
3. \dotfill
}
\end{tcolorbox}

\vspace{4pt}

% Grille d'auto-évaluation
\begin{tcolorbox}[
  enhanced,
  colback=neutrallight,
  colframe=secondary!60!cardborder,
  arc=2mm,boxrule=0.7pt,
  top=4pt,bottom=4pt,left=7pt,right=7pt,
  title={\small\textbf{\faTasks\ Grille d'auto-évaluation de l'élève}},
  coltitle=primary,colbacktitle=secondarylight,
  titlerule=0.5pt
]
{\footnotesize
\begin{tabularx}{\linewidth}{Xcc}
\textbf{Critère de réussite} & \textbf{Oui \faSmile} & \textbf{Non \faFrown} \\
\midrule
J'ai bien identifié la règle grammaticale. & \qcmCircle & \qcmCircle \\
J'ai réussi les exercices d'application sans aide. & \qcmCircle & \qcmCircle \\
J'ai écrit des phrases correctes et soignées. & \qcmCircle & \qcmCircle \\
\end{tabularx}
}
\end{tcolorbox}

\end{document}
